From \parencite{adekoya2022applications_ai_em} :

“Several diverse theories, algorithms, and methods are available, which presents a challenge and a barrier in choosing the right AI technique for the appropriate engineering process or manufacturing process and environments.”

“Notwithstanding the benefits of AI techniques in engineering and manufacturing, Wuest et al. (2016) argue that its diversity poses a challenge to countless manufacturing practitioners and engineers regarding the proper technique to select from among the lots.”

“Therefore, it might hinder the utilisation of AI techniques by an appreciable percentage of experts in this field.”

“Hence, there are no fast rules of superiority among ML algorithms.”

“Thus, their performance may differ from one area of application to the other.”

“Besides, the pertinent literature is disseminated over various journals, conference proceedings, and research communities.”

“Hence, conducting a systematic survey to scrutinise and classify the existing literature is worthwhile.”

“However, most of these existing systematic review studies surveyed the applications of a single AI technique [...] Otherwise, cover only one engineering or manufacturing domain [...].”

“Thus, the progress and advancement of AI techniques (like ML and DL) in various engineering and manufacturing systems and environments have not yet been well addressed in the literature.”


From \parencite{wuest2016ml_manufacturing} :

“One area, which saw fast pace developments in terms of not only promising results but also usability, is machine learning.”

“Promising an answer to many of the old and new challenges of manufacturing, machine learning is widely discussed by researchers and practitioners alike.”

“The application of ML techniques increased over the last two decades due to various factors, e.g. the availability of large amounts of complex data [...] and the increased usability and power of available ML tools.”

“ML is already widely applied in different areas of manufacturing, e.g. optimization, control, and troubleshooting.”

“However, the field is very broad and even confusing which presents a challenge and a barrier hindering wide application.”

“However, the field of machine learning is very diverse and many different algorithms, theories, and methods are available.”

“For many manufacturing practitioners, this represents a barrier regarding the adoption of these powerful tools.”

“A major challenge of increasing importance is the question what ML technique and algorithm to choose.”

“There are many different ML methods, tools, and techniques available, each with distinct advantages and disadvantages.”

“The best fitting algorithm has to be found in testing various ones in a realistic environment.”


From \parencite{plathottam2023review_ai_manufacturing} :

“The predictive and analytic capabilities of AI/ML models can provide useful insights to human operators to facilitate planning, support real-time decision-making, and improve manufacturing efficiency and safety.”

“The increasing diversity and complexity of manufacturing tasks, workflows, and supply chains has increased the number of variables and interdependencies.”

“Implementing AI solutions, including using mature AI/ML technologies remains challenging.”

“the choice of which AI/ML solutions to implement for specific manufacturing problems is not always trivial and involves a certain degree of trade-offs.”

“AI/ML implementation decisions must be suited to each company's unique situation and needs”

“pressure from within the industry to use AI has contributed to companies feeling obligated to pursue AI/ML solutions even when they lack a concrete plan.”


From \parencite{weichert2019review_ml_optimization_production}

“there is tremendous progress and large interest in integrating machine learning and optimization methods on the shop floor in order to improve production processes.”

“machine learning [is used] to save energy, time, and resources, and avoid waste.”

“The interest in the integration of machine learning and optimization algorithms in production processes has increased steadily since 2009 and has been at a consistently high level since 2014.”

“The first step to solving a machine learning problem is to identify the data that is available.”

“The knowledge of the structure of production data is crucial for the right choice of machine learning models.”

“The reviewed literature shows that the production-related applications of machine learning with or without optimization are manifold.”

“there is hardly any correlation between the used data, the amount of data, the machine learning algorithms, the used optimizers, and the respective problem from the production.”

“process complexity, the number of data points per input dimension to the model, and the model complexity are highly related.”

“the connection between the process complexity, the stored data amount, and the model complexity [...] was not respected”

“this connection was not respected, leading to complex models being trained on low amounts of data, risking overfitting and/or a lack of interpretability.”


From \parencite{rai2021ml_manufacturing_industry4} :

“The machine learning (ML) field has deeply impacted the manufacturing industry in the context of the Industry 4.0 paradigm.”

“ML techniques enable the generation of actionable intelligence by processing the collected data to increase manufacturing efficiency without significantly changing the required resources.”

“the ability of ML techniques to provide predictive insights has enabled discerning complex manufacturing patterns and offers a pathway for an intelligent decision support system”

“Advances in Big data and Machine Learning (ML) is changing the traditional manufacturing era into the smart manufacturing era of Industry 4.0”

“ML techniques [...] has the potential to become the main driver in uncovering fine-grained intricate production patterns in smart manufacturing paradigm and offering timely decision support”

“While different ML techniques have been used in a variety of manufacturing applications in the past, many open questions and challenges remain”

“Manufacturing is undertaking a definitive shift as it assimilates and is transformed by the usage of machine learning techniques.”

“The utility of ML in manufacturing applications reflects that ML can be integrated into every phase of product lifecycle, starting from design to disposal of the product in a manufacturing environment.”


From \parencite{chen2023ml_manufacturing_industry4} :

“Machine Learning (ML) is playing a critical role in the digitalization of manufacturing operations towards Industry 4.0.”

“Recently, ML has been applied in several fields of production engineering to solve a variety of tasks with different levels of complexity and performance.”

“However, in spite of the enormous number of ML use cases, there is no guidance or standard for developing ML solutions from ideation to deployment.”

“it remains very difficult for non-experts working in the manufacturing industry to begin developing ML solutions for their specific problems.”

“The first challenging part of application is to formulate the actual problems to be solved.”

“the implementation of ML solution for different tasks in an industrial environment from scratch has not yet been fully discussed.”

“A large number of ML use cases have shown the great potential for addressing complex manufacturing problems”

“there are critical challenges to overcome before the successful application of ML in manufacturing can be realized.”

“To fill the gap between academic research and manufacturing industries, we provided an actionable pipeline for the implementation of ML solutions by production engineers from ideation through to deployment”

“We hope that this can provide support in method selection for decision-makers considering ML solutions.”


From \parencite{sonntag2023pattern_ml_product_development} :

“companies lack the implementation of these approaches… due to inadequate knowledge and experience”

“the initial guess of suitable algorithms requires a certain amount of experience.”

“only a fraction of them actually do so”

“75% of the companies surveyed believe that automation approaches… have no significance for them”

“only 16% plan to implement automation approaches”

“lack of experience and knowledge… to recognize opportunities… and identify suitable implementation options”

“To use ML… one has to identify the most suitable ML algorithm out of a large number”

“the procedure is based on… trial-and-error”

“this is… laborious”

“In summary, no solution allows for an informed selection of ML algorithms based on a well-defined
problem within the product development process."

“support non-experts in identifying the optimal algorithm”

“make the knowledge… accessible for non-experts”

“pattern language… enabling problem solving by non-experts”

“patterns… make expert knowledge easily accessible”

From \parencite{arena2024conceptual_framework_ml_algorithm_selection_predictive_maintenance} :

“the need to choose the right algorithm for a specific task arises as a challenging issue”

“there is no unique ML algorithm that can efficiently handle the analysis of every dataset or all use cases”

“this may frequently lead to results that are poor or too generic for the problem at hand”

“it is crucial to define the applicability domain”

“no method is better than all the others for all possible datasets”

“an important step in developing an ML model is deciding which learning method provides the best results”

“most of the articles focus on the development of the ML model… not considering what conditions lead to selecting one rather than another”

“many of them focus on accuracy… prediction speed… or computational requirements”

“no tool relates algorithmic approaches to the maintenance process”

“theoretical accuracy and the computational speed do not ensure satisfactory results”

“making the algorithm selection a non-trivial task”

“implementing an ML-based PdM strategy presents multiple challenges”

From \parencite{presciuttini2024ml_iot_manufacturing_operations} :

“The lack of interpretability of most ML and DL applications is one of the key barriers to their adoption in real industrial environments”

“it is crucial to go beyond their ‘black box’ nature by emphasizing their interpretability dimension”

“most of the works are highly focused on the predictive power of the developed models, without considering the reasons leading to the obtained output”

“AI methodologies are often regarded as black boxes, with little or no effort made to clarify the outcomes”

“it is crucial to provide an explanation of why the solution works and interpret the results for widespread adoption in manufacturing”

“While limiting ML and DL models to predictive purposes aids in issue anticipation, it falls short of identifying the actual causes of inefficiencies”

“interpretable models can help maintenance teams understand the factors that mostly contribute to equipment failures”

“the trade-off between model interpretability and accuracy could be explored”

From \parencite{paleyes2022challenges_deploying_ml_survey_case_studies} :

“there are significant differences between what works in an academic setting and what is required by a real world system”

“the question becomes not ‘Where is it used?,’ but rather ‘How difficult is it to use?’”

“practitioners face issues at each stage of the deployment process”

“the overall effectiveness of the solution depends on the training and test data as much as on the algorithm”

“in many practical cases the selection of a model is decided by one key characteristic of a model: complexity”

“in practice simpler models are often chosen”

“the ability to interpret the output of a model… often plays a critical role in model selection and can even outweigh performance considerations”

“an increase in model performance does not translate into a gain in business value”

“hyper-parameter optimization… is very expensive and resource-heavy in practice”

“it is important to understand the most critical pain points and to provide tools, services, and best practices that address those points”

From \parencite{ali2017accurate_multi_criteria_decision_making_methodology_recommending_machine_learning_algorithm} :

“Manual evaluation of machine learning algorithms and selection of a suitable classifier… is highly time consuming and challenging task.”

“If the selection is not carefully and accurately done, the resulting classification model will not be able to produce the expected performance results.”

“The heuristics-based algorithm(s) selection is a risky task and sometimes result in selection of a sub-optimal performance algorithm(s).”

“The datasets have different intrinsic characteristics, and the candidate classifiers have different capabilities and strengths.”

“This process become more challenging when the selection of best classifier is based on multiple-criteria under strict conditions and constraints.”

“No machine learning algorithm performs well on all kind of learning problems.”

“The evaluation only on the basis of accuracy may misleads the selection of optimum performance algorithm.”

“The objective of multi-criteria evaluation is to balance the trade-off between these criteria rather than maximizing a single criterion.”

“The main issue in multi-criteria evaluation is the selection and prioritization of suitable criteria.”

“No classification algorithms is superior on all problems… if one algorithm outperforms others on one criterion, it may underperforms on other criteria.”

From \parencite{dogan2021machine_learning_data_mining_manufacturing} :

“ML and data mining (DM) applications came to existence in this sector nearly two decades ago to solve manufacturing problems.”

“The machine learning field [...] has been considered one of the most promising improvements in the manufacturing domain.”

“Modern manufacturing plants use powerful data acquisition systems to electronically collect and transfer data from almost all the processes of the organization.”

“A very common challenge of ML application in manufacturing is the pre-processing of data since it has a critical impact on the results.”

“Another key challenge is the question of what the ML method and algorithm to choose.”

“The success of each ML technique depends on the data structure on which it is performed.”

“However, there are no standardized rules that indicate which data preprocessing techniques should be applied for a specific type of manufacturing problem.”

“Every knowledge discovery operation is unique in its context. Therefore, the optimal ML technique which is capable of solving the targeted manufacturing problem is not clear previously.”

“Experts who have some experiences can only propose some kind of approach by referring to past studies. However, those suggestions do not guarantee positive results for that context.”

“Each manufacturing problem is different and the performance of each algorithm depends on the data available, as well as parameter settings.”

"Machine learning applications will probably keep growing at a higher rate particularly in manufacturing because computing power is increasing day by day and the size of available data is much greater than it was in a few years ago."

"Every knowledge discovery operation is unique in its context. Therefore, the optimal ML technique which is capable of solving the targeted manufacturing problem is not clear previously."

"Each manufacturing problem is different and the performance of each algorithm depends on the data available, as well as parameter settings."

"since manufacturers have solely worked in the manufacturing sector, they do not have strong background information and experience about ML methods and their usage."

"Experts who have some experiences can only propose some kind of approach by referring to past studies. However, those suggestions do not guarantee positive results for that context."

Alternative algorithms and parameter trials bring about time usage costs. For this reason, it has taken too much time to put machine learning methods in daily practice in the manufacturing sector."

"ML and DM applications came to existence in this sector nearly two decades ago to solve manufacturing problems."

From \parencite{sala2018selecting_ml_algorithm} :

“The considerable number of works, depicting a wide variety of algorithms and widespread applications, creates an extensive knowledge on the topic. At the same time, it may also generate disorientation in the selection of the right approach.”

“Thus, the need of synthesis and guidelines to drive the selection of the most suitable algorithm for a specific scope arises.”

“However, a unique ML algorithm able to deal with the analysis of every dataset does not exist, making the selection of the right ML algorithm challenging.”

“Some of the frameworks currently available tend to use as main selection drivers the accuracy and/or prediction speed of the ML algorithms.”

“The problem in this class of selection frameworks is that they may result too general most of the time.”

“In this way, unsuitable ML algorithms may be selected leading in turn to poor results.”

“The problem with these criteria is that they require that the user knows the characteristics of each model and that s/he computes the value for each single model under consideration.”

“In this case, the model selection could be biased by the initial pool of algorithms selected by the user that could not include the most suitable ones.”

“In the selection of a suitable ML algorithm for data analysis many aspects should be taken into account.”

“Most of the times selecting an algorithm only on the base of the promised accuracy or computational speed leads to unsatisfactory results.”

From \parencite{lee2020machine_learning_enterprises} :

“With a wide variety of machine-learning algorithms available, managers would find it challenging to select the best ones to solve their problems.”

“Selecting the best machine-learning algorithm for the data set is a challenge for any machine-learning project.”

“No one algorithm works best for all problems.”

“Several variables, such as the volume, variety, and velocity of data, as well as the type of problem at hand, may affect the performance of machine-learning algorithms.”

“Balancing the trade-off between accuracy and interpretability can be challenging.”

“Managers should ideally evaluate multiple machine-learning algorithms in parallel to determine the highest achievable accuracy.”

From \parencite{wolpert1997no_free_lunch} :

For any algorithm, any elevated performance over one class of problems is offset by performance over another class."

"If some algorithm's performance is superior to that of another algorithm over some set of optimization problems, then the reverse must be true over the set of all other optimization problems."

"If knowledge of problem characteristics is not incorporated into the optimization algorithm, then there are no formal assurances that the algorithm chosen will be at all effective."

From \parencite{smith_miles2009algorithm_selection} :

"The Algorithm Selection Problem seeks to answer the question: which algorithm is likely to perform best for my problem?"

"From the early work of John Rice which formalized the Algorithm Selection Problem, to the No Free Lunch Theorems of Wolpert and Macready which stated that for any algorithm, any elevated performance over one class of problems is exactly paid for in performance over another class, it has long been acknowledged that there is no universally-best algorithm to tackle a broad problem domain."

"The NFL theorems suggest that we need to move away from a 'black-box' approach to algorithm selection. We need to understand more about the characteristics of the problem in order to select the most appropriate algorithm."